\documentclass[10pt,a4paper]{amsart}
\usepackage{amsmath,amssymb,amsthm,graphicx}
% Retained preamble for the Stacks Project chapter snapshot in this directory.
% The chapter's own preamble.tex is a copy of the shared Stacks Project preamble; the block
% below is sliced MECHANICALLY, line for line, from it. Nothing was added, removed or
% redefined; the wrapper class, the xr-hyper/hyperref wiring and the externaldocument list
% are omitted because the assembler supplies the document class.
%
% lines taken from preamble.tex: 3, 5, 9, 11, 12, 156, 157, 158, ... (51 lines total)
%
% package coverage restored from the chapter's own preamble (2026-10-13): xy, xy-2cell

\usepackage[all]{xy}
\xyoption{2cell}
\UseAllTwocells

\newenvironment{reference}{\comment}{\endcomment}
\newenvironment{slogan}{\comment}{\endcomment}
\newenvironment{history}{\comment}{\endcomment}
\theoremstyle{plain}
\newtheorem{theorem}[subsection]{Theorem}
\newtheorem{proposition}[subsection]{Proposition}
\newtheorem{lemma}[subsection]{Lemma}
\theoremstyle{definition}
\newtheorem{definition}[subsection]{Definition}
\newtheorem{example}[subsection]{Example}
\newtheorem{exercise}[subsection]{Exercise}
\newtheorem{situation}[subsection]{Situation}
\theoremstyle{remark}
\newtheorem{remark}[subsection]{Remark}
\newtheorem{remarks}[subsection]{Remarks}
\numberwithin{equation}{subsection}
\def\lim{\mathop{\mathrm{lim}}\nolimits}
\def\colim{\mathop{\mathrm{colim}}\nolimits}
\def\Spec{\mathop{\mathrm{Spec}}}
\def\Hom{\mathop{\mathrm{Hom}}\nolimits}
\def\Ext{\mathop{\mathrm{Ext}}\nolimits}
\def\SheafHom{\mathop{\mathcal{H}\!\mathit{om}}\nolimits}
\def\SheafExt{\mathop{\mathcal{E}\!\mathit{xt}}\nolimits}
\def\Sch{\mathit{Sch}}
\def\Mor{\mathop{\mathrm{Mor}}\nolimits}
\def\Ob{\mathop{\mathrm{Ob}}\nolimits}
\def\Sh{\mathop{\mathit{Sh}}\nolimits}
\def\NL{\mathop{N\!L}\nolimits}
\def\CH{\mathop{\mathrm{CH}}\nolimits}
\def\proetale{{pro\text{-}\acute{e}tale}}
\def\etale{{\acute{e}tale}}
\def\QCoh{\mathit{QCoh}}
\def\Ker{\mathop{\mathrm{Ker}}}
\def\Im{\mathop{\mathrm{Im}}}
\def\Coker{\mathop{\mathrm{Coker}}}
\def\Coim{\mathop{\mathrm{Coim}}}
\DeclareMathSymbol{\boxtimes}{\mathbin}{AMSa}{"02}
\def\QCohstack{\mathcal{QC}\!\mathit{oh}}
\def\Cohstack{\mathcal{C}\!\mathit{oh}}
\def\Spacesstack{\mathcal{S}\!\mathit{paces}}
\def\Quotfunctor{\mathrm{Quot}}
\def\Hilbfunctor{\mathrm{Hilb}}
\def\Curvesstack{\mathcal{C}\!\mathit{urves}}
\def\Polarizedstack{\mathcal{P}\!\mathit{olarized}}
\def\Complexesstack{\mathcal{C}\!\mathit{omplexes}}
\def\Pic{\mathop{\mathrm{Pic}}\nolimits}
\def\Picardstack{\mathcal{P}\!\mathit{ic}}
\def\Picardfunctor{\mathrm{Pic}}
\def\Deformationcategory{\mathcal{D}\!\mathit{ef}}
% lane-s11: discard bibliographic \begin{reference} blocks entirely (incl. any \cite inside).
\newsavebox{\limrefbox}
\ifcsname reference\endcsname\expandafter\let\csname reference\endcsname\undefined\fi\ifcsname endreference\endcsname\expandafter\let\csname endreference\endcsname\undefined\fi
\def\citation#1{\relax}
\newenvironment{reference}{\setbox\limrefbox=\vbox\bgroup\let\cite\relax}{\egroup\ignorespaces}
% lane-s11: the chapter's retained preamble defines slogan/history via comment.sty's
% \comment macro, which is never loaded; redefine both with the same guarded
% discard pattern (content is authorial commentary, kept verbatim in the spans).
\ifcsname slogan\endcsname\expandafter\let\csname slogan\endcsname\undefined\fi\ifcsname endslogan\endcsname\expandafter\let\csname endslogan\endcsname\undefined\fi
\newenvironment{slogan}{\setbox\limrefbox=\vbox\bgroup}{\egroup\ignorespaces}
\ifcsname history\endcsname\expandafter\let\csname history\endcsname\undefined\fi\ifcsname endhistory\endcsname\expandafter\let\csname endhistory\endcsname\undefined\fi
\newenvironment{history}{\setbox\limrefbox=\vbox\bgroup}{\egroup\ignorespaces}
\usepackage{cleveref}
\crefname{originalreference}{original reference}{original references}
\Crefname{originalreference}{Original reference}{Original references}
\makeatletter
\newlabel{morphisms-lemma-integral-universally-closed}{{01WM}{source}{}{sourcepointer-1}{}}
\newlabel{morphisms-lemma-integral-universally-closed@cref}{{[originalreference][0][]01WM}{[1][0][]source}{}{}{}}
\makeatother
\crefname{originalreference}{original reference}{original references}
\Crefname{originalreference}{Original reference}{Original references}
\makeatletter
\newlabel{morphisms-lemma-image-universally-closed-separated}{{09MQ}{source}{}{sourcepointer-2}{}}
\newlabel{morphisms-lemma-image-universally-closed-separated@cref}{{[originalreference][0][]09MQ}{[1][0][]source}{}{}{}}
\makeatother
\crefname{originallemma}{original lemma}{original lemmas}
\Crefname{originallemma}{Original lemma}{Original lemmas}
\makeatletter
\newlabel{lemma-finite-in-finite-and-finite-presentation}{{09YY}{source}{}{sourcepointer-3}{}}
\newlabel{lemma-finite-in-finite-and-finite-presentation@cref}{{[originallemma][0][]09YY}{[1][0][]source}{}{}{}}
\makeatother
\crefname{originallemma}{original lemma}{original lemmas}
\Crefname{originallemma}{Original lemma}{Original lemmas}
\makeatletter
\newlabel{lemma-limit-affine}{{01Z6}{source}{}{sourcepointer-4}{}}
\newlabel{lemma-limit-affine@cref}{{[originallemma][0][]01Z6}{[1][0][]source}{}{}{}}
\makeatother
\crefname{originalproposition}{original proposition}{original propositions}
\Crefname{originalproposition}{Original proposition}{Original propositions}
\makeatletter
\newlabel{proposition-approximate}{{01ZA}{source}{}{sourcepointer-5}{}}
\newlabel{proposition-approximate@cref}{{[originalproposition][0][]01ZA}{[1][0][]source}{}{}{}}
\makeatother
\crefname{originallemma}{original lemma}{original lemmas}
\Crefname{originallemma}{Original lemma}{Original lemmas}
\makeatletter
\newlabel{lemma-descend-finite-presentation}{{01ZM}{source}{}{sourcepointer-6}{}}
\newlabel{lemma-descend-finite-presentation@cref}{{[originallemma][0][]01ZM}{[1][0][]source}{}{}{}}
\makeatother
\crefname{originallemma}{original lemma}{original lemmas}
\Crefname{originallemma}{Original lemma}{Original lemmas}
\makeatletter
\newlabel{lemma-descend-finite-finite-presentation}{{01ZO}{source}{}{sourcepointer-7}{}}
\newlabel{lemma-descend-finite-finite-presentation@cref}{{[originallemma][0][]01ZO}{[1][0][]source}{}{}{}}
\makeatother
\crefname{originalreference}{original reference}{original references}
\Crefname{originalreference}{Original reference}{Original references}
\makeatletter
\newlabel{coherent-lemma-image-affine-finite-morphism-affine-Noetherian}{{01YQ}{source}{}{sourcepointer-8}{}}
\newlabel{coherent-lemma-image-affine-finite-morphism-affine-Noetherian@cref}{{[originalreference][0][]01YQ}{[1][0][]source}{}{}{}}
\makeatother
\title{Affineness of the limit in the separated finite presentation case}
\author{The Stacks Project Authors}
\date{stacks-project:limits.tex:local-snapshot-20261012; GNU FDL 1.2 or later}
\begin{document}
\maketitle
\noindent\textit{Editorial note.} No omissions, hints or 'immediate' steps in the counted span.\par\medskip
\noindent\textit{Allowed standard prerequisites (editorial).}\begin{itemize}
\item Affine detection. Limits, affine detection lemmas of the chapter
\end{itemize}\medskip
\ifdefined\hypertarget\hypertarget{sourcepointer-1}{}\fi
\noindent\textit{Source reference.} Original reference 01WM: sources/stacks-chow/tags.txt, source lines 2487--2487. Cross-chapter standard prerequisite cited by the author inside the counted span; it is used as a pointer and its own proof is not part of this item.\par\medskip
\ifdefined\hypertarget\hypertarget{sourcepointer-2}{}\fi
\noindent\textit{Source reference.} Original reference 09MQ: sources/stacks-chow/tags.txt, source lines 12188--12188. Cross-chapter standard prerequisite cited by the author inside the counted span; it is used as a pointer and its own proof is not part of this item.\par\medskip
\ifdefined\hypertarget\hypertarget{sourcepointer-3}{}\fi
\noindent\textit{Source reference.} Original lemma 09YY: sources/stacks-chow/tags.txt, source lines 12581--12581. Cross-chapter standard prerequisite cited by the author inside the counted span; it is used as a pointer and its own proof is not part of this item.\par\medskip
\ifdefined\hypertarget\hypertarget{sourcepointer-4}{}\fi
\noindent\textit{Source reference.} Original lemma 01Z6: sources/stacks-chow/tags.txt, source lines 2581--2581. Cross-chapter standard prerequisite cited by the author inside the counted span; it is used as a pointer and its own proof is not part of this item.\par\medskip
\ifdefined\hypertarget\hypertarget{sourcepointer-5}{}\fi
\noindent\textit{Source reference.} Original proposition 01ZA: sources/stacks-chow/tags.txt, source lines 2585--2585. Cross-chapter standard prerequisite cited by the author inside the counted span; it is used as a pointer and its own proof is not part of this item.\par\medskip
\ifdefined\hypertarget\hypertarget{sourcepointer-6}{}\fi
\noindent\textit{Source reference.} Original lemma 01ZM: sources/stacks-chow/tags.txt, source lines 2599--2599. Cross-chapter standard prerequisite cited by the author inside the counted span; it is used as a pointer and its own proof is not part of this item.\par\medskip
\ifdefined\hypertarget\hypertarget{sourcepointer-7}{}\fi
\noindent\textit{Source reference.} Original lemma 01ZO: sources/stacks-chow/tags.txt, source lines 2601--2601. Cross-chapter standard prerequisite cited by the author inside the counted span; it is used as a pointer and its own proof is not part of this item.\par\medskip
\ifdefined\hypertarget\hypertarget{sourcepointer-8}{}\fi
\noindent\textit{Source reference.} Original reference 01YQ: sources/stacks-chow/tags.txt, source lines 2565--2565. Cross-chapter standard prerequisite cited by the author inside the counted span; it is used as a pointer and its own proof is not part of this item.\par\medskip
\begin{lemma}
\label{lemma-affine}
\begin{slogan}
A scheme, admitting a finite surjective map from an affine scheme, is affine.
\end{slogan}
Let $f : X \to S$ be a morphism of schemes.
Assume that $f$ is surjective and finite, and assume that $X$ is affine.
Then $S$ is affine.
\end{lemma}

\begin{proof}
Since $f$ is surjective and $X$ is quasi-compact we see that $S$ is
quasi-compact. Since $X$ is separated and $f$ is surjective and
universally closed (Morphisms, Lemma
\ref{morphisms-lemma-integral-universally-closed}), we see that $S$
is separated (Morphisms, Lemma
\ref{morphisms-lemma-image-universally-closed-separated}).

\medskip\noindent
By Lemma \ref{lemma-finite-in-finite-and-finite-presentation}
we can write $X = \lim_a X_a$ with $X_a \to S$ finite and of finite
presentation. By Lemma \ref{lemma-limit-affine} we see that $X_a$
is affine for some $a \in A$. Replacing $X$ by $X_a$ we may assume
that $X \to S$ is surjective, finite, of finite presentation and
that $X$ is affine.

\medskip\noindent
By Proposition \ref{proposition-approximate} we may write
$S = \lim_{i \in I} S_i$ as
a directed limits of schemes of finite type over $\mathbf{Z}$.
By Lemma \ref{lemma-descend-finite-presentation} we can
after shrinking $I$ assume there exist schemes $X_i \to S_i$
of finite presentation such that $X_{i'} = X_i \times_S S_{i'}$
for $i' \geq i$ and such that $X = \lim_i X_i$. By
Lemma \ref{lemma-descend-finite-finite-presentation} we may
assume that $X_i \to S_i$ is finite for all $i \in I$ as well.
By Lemma \ref{lemma-limit-affine} once again we may assume that $X_i$ is
affine for all $i \in I$. Hence the result follows from the
Noetherian case, see Cohomology of Schemes,
Lemma \ref{coherent-lemma-image-affine-finite-morphism-affine-Noetherian}.
\end{proof}

\end{document}
